\scanpage{11}
Réciproquement, partons d'une équation comme dessus, où les \deleted{$a$} \inserted{coeff. $a,b$ ---} sont des sections de $\underline O_S$. On sait qu'\deleted{\unsure{elle|ce}} \inserted{une telle équation} définit un diviseur relatif \inserted{positif} de $\mathbb P^2_S/S$, soit $C$. Il reste à écrire que c'est là une \uline{courbe simple} sur $S$. On sait qu'il f et s pour celà que le discriminant de la forme
\[
\delta(a,b,c,d,e,f)
\]
soit une \uline{unité} partout. Ici le discriminant s'explicite : ---

On a obtenu en résumé

\uline{Proposition} La donnée d'une courbe elliptique sur $S$, \deleted{\unsure{avec}} \inserted{munie} d'une \deleted{``base adaptée''} \inserted{``base spéciale''} de $\mathcal E_3(C/S)$, \deleted{``\unsure{adaptée}''} est équival\deleted{\unsure{ents}}ente à la donnée de sections $a,b,c,d,e,f$ de $\underline O_S$ satisfaisant à $\delta(a,b,c\unsure{,|+}d,e,f)=0$, cette correspondance est compatible avec l'image inverse.

On déduit de ceci : Considérons le schéma affine de type fini sur $\mathbb Z$
\[
M=\mathrm{Spec}\bigl(\mathbb Z[A,B,C,D,E,F][\delta^{-1}]\bigr)
\]
qui est un ouvert affine de $\mathrm{Spec}\,\mathbb Z[A,B,C,D,E,F]\simeq\mathbb E^6_{\mathbb Z}$. \deleted{Faisons opérer \unsure{dans ce|dessus}}
% archive sheet number: 10
